\documentclass[11pt]{article}
\usepackage[margin=1in]{geometry}
\usepackage{hyperref}
\usepackage{listings}
\usepackage{xcolor}
\usepackage{amsmath}

\lstset{
  basicstyle=\ttfamily\small,
  breaklines=true,
  columns=fullflexible,
  frame=single,
  backgroundcolor=\color{gray!10},
  showstringspaces=false
}

\title{URI OPeNDAP Test Case: Working Log}
\author{Compiled for P. Cornillon}
\date{2026-05-16}

\begin{document}
\maketitle

\section*{Purpose}

This file records commands, intermediate observations, and decisions
made during the exploratory analysis of the URI OPeNDAP server at
\url{https://sst-aqua.gso.uri.edu/opendap/}. Findings are written to
\texttt{docs/uri\_test\_case.tex}; what lives here is the
\emph{process}, including reasoning that did not make it into the
findings document.

Each section is timestamped. The session container reports UTC.

\section{Prompt \#1 (2026-05-16 14:01 UTC)}

\subsection{Instruction}

\emph{``Read this document. Explore the contents of the URI OPeNDAP
server. Describe the data and the metadata. Describe the data within
the server and explain your reasoning. Be mindful of the overall
project that we are working on. Spend up to 1 hour on this task.''}

\subsection{Plan}

\begin{enumerate}
  \item Confirm the server is reachable. Identify the server software
    (Hyrax vs.\ THREDDS vs.\ ERDDAP) and the protocol version (DAP2 /
    DAP4) from the landing page and from probing
    \texttt{.dmr}/\texttt{.dds} on a sample dataset.
  \item Walk the top-level catalog. Note container names and
    file-name patterns~--- since the project frames \texttt{grep-dap}
    around inferring contents in the face of loose semantic
    metadata, directory and file names are first-order signal.
  \item For each \emph{publicly accessible} top-level container,
    pick a representative dataset and pull its metadata (DMR if
    DAP4, DDS+DAS otherwise). Record dimensions, shapes, dtypes,
    attributes.
  \item Pull small slices with \texttt{pydap} to sanity-check the
    metadata (e.g.\ that a variable claiming Kelvin really sits near
    275--310, that lat/lon really span $\pm$90 / $\pm$180).
  \item Synthesize findings with explicit caveats about what is
    inferred vs.\ stated.
\end{enumerate}

\subsection{Reachability and server identification}

The landing page returns HTTP\,200 from \texttt{nginx}, and the body is
a Hyrax-rendered \texttt{contents.html} (CSS \texttt{contents.css},
logo \texttt{logo.png}, DAP4 + DAP2 link columns). So this is a
\textbf{Hyrax} server. Both DAP4 and DAP2 endpoints are advertised on
each leaf. The reverse-proxy URL exposes an internal port
\texttt{:1512}, which suggests \texttt{nginx} terminating TLS in front
of the Hyrax / BES backend.

\begin{lstlisting}
$ curl -sI https://sst-aqua.gso.uri.edu/opendap/
HTTP/2 200
server: nginx
content-type: text/html
content-description: dap_directory
\end{lstlisting}

\subsection{Top-level catalog walk}

Seventeen entries at the root. Public reads succeed on two of them
(\texttt{SST\_Orbits/}, \texttt{gradients\_by\_period/}); the rest
return Hyrax 403 ``\emph{Not allowed to access
/usr/share/hyrax/$\langle$dir$\rangle$/}'' --- i.e.\ filesystem
permissions on the BES side, not an authentication challenge. The 403
inventory tells us a lot anyway, because the names follow the URI
group's well-known naming conventions.

\begin{lstlisting}
ROOT (HTTP code observed for contents.html):
  @eaDir/                            -  (Synology NAS housekeeping, ignore)
  JAXA_Orbits/                       403
  JAXA_Orbits_OLD/                   403
  MUR/                               403
  RSS_Orbits/                        403
  RSS_Orbits_OLD/                    403
  RSS_Orbits_REALLY_OLD/             403
  SST_Orbits/                        200  <-- public
  gradients_by_period/               200  <-- public
  gradients_by_period_5day/          403
  iQuamBuoy/                         403
  matchups/                          403
  matchups_v2/                       403
  matchups_v3/                       403
  timeSeries/                        403
  timeSeries_v2/                     403
  timeSeries_v2_backup/              403
\end{lstlisting}

The two public branches behave very differently:
\begin{itemize}
  \item \texttt{SST\_Orbits/} is partitioned by year then month
    (\texttt{2002}--\texttt{2024}, then \texttt{01}--\texttt{12}); the
    leaves are NetCDF-4 files whose names encode the satellite, the
    orbit number, the start time, and the processing tag:
    \begin{quote}\ttfamily AQUA\_MODIS\_orbit\_115220\_\\20240101T011558\_L2\_SST-URI\_24-2.nc4\end{quote}
    Each year contains 12 month subdirs; each month contains roughly
    14 orbits/day $\times$ 28--31 days $\approx$ 400--430 files.
    Approximate inventory: $\sim$23 years $\times$ 14.4 orbits/day
    $\times$ 365 d $\approx$ \textbf{$\sim 1.2 \times 10^5$ orbit
    files}.
  \item \texttt{gradients\_by\_period/} is flat: 270 NetCDF files
    \texttt{augmented\_monthly\_stats\_MM\_YYYY.nc} spanning
    \texttt{01\_2003} through \texttt{12\_2024}, with a small surplus
    over $22 \times 12 = 264$ that I have not chased down.
\end{itemize}

\subsection{Metadata fetch: SST\_Orbits sample}

Picked the first orbit of 2024 (orbit 115220, $\sim$01:15:58 UTC,
2024-01-01). DMR XML is 141\,lines of structure plus another
$\sim$900\,lines of attributes; clearly DAP4.

\begin{lstlisting}
$ curl -s "$URL.dmr.xml" | head -120
<?xml version="1.0" encoding="ISO-8859-1"?>
<Dataset ... dapVersion="4.0" ...
    <Dimension name="ny" size="40271"/>
    <Dimension name="nx" size="1354"/>
    <Dimension name="i"  size="5"/>
    ...
\end{lstlisting}

Root variables (selected, after parsing the DMR with a small
\texttt{re}-based script in \texttt{uri\_walk\_public.py}):

\begin{itemize}
  \item \texttt{SST\_In} (Int16, ny$\times$nx, scale\_factor 0.005,
    units \texttt{C}, standard\_name
    \texttt{sea\_surface\_temperature}).
  \item \texttt{regridded\_sst} (Int16, ny$\times$nx, units \texttt{C},
    standard\_name \texttt{masked\_sea\_surface\_temperature}).
  \item \texttt{eastward\_gradient}, \texttt{northward\_gradient}
    (Int32, ny$\times$nx, scale 0.0001, units \texttt{C/km},
    standard\_name \texttt{eastward\_temperature\_gradient} /
    \texttt{northward\_temperature\_gradient}).
  \item \texttt{refined\_mask}, \texttt{qual\_sst} (Int8 quality /
    flag arrays).
  \item \texttt{latitude}, \texttt{longitude} (Int32 ny$\times$nx) ---
    pixel locations. Note: the DMR does \emph{not} declare a
    scale\_factor for these. Comparison with the Float32
    \texttt{nadir\_latitude/longitude} (which is in degrees) shows
    the Int32 lat/lon values are 1000$\times$ the degree value; the
    implicit scale is 0.001. That is a metadata gap a client must
    repair by inference.
  \item \texttt{regridded\_latitude}, \texttt{regridded\_longitude}
    --- same caveat.
  \item Swath geometry: \texttt{left/right/nadir\_latitude},
    \texttt{left/right/nadir\_longitude} (Float32, ny;
    degrees\_north / degrees\_east).
  \item \texttt{time\_from\_start\_orbit} (Float32, ny, seconds).
  \item \texttt{DateTime} (Float64 scalar, seconds since 1970-01-01;
    standard\_name \texttt{time}). For the sample file the value is
    \texttt{1704071758.2 s} $\to$ 2024-01-01T01:15:58 UTC, matching
    the filename exactly.
  \item \texttt{region\_start}, \texttt{region\_end} (Int32, i=5) ---
    presumably five along-orbit subregions (mid- vs.\ polar bands?).
\end{itemize}

The orbit file also carries two \texttt{Group}s:
\begin{itemize}
  \item \texttt{/Regrid\_to\_L2eqa/} --- companion fields regridded
    to an L2 equal-area grid (its own \texttt{nxL2eqa},
    \texttt{nyL2eqa}, plus \texttt{nxAMSR\_E}, \texttt{nyAMSR\_E}
    dimensions). Variables include
    \texttt{L2eqa\_MODIS\_SST}, \texttt{L2eqa\_MODIS\_std\_SST},
    \texttt{L2eqa\_MODIS\_num\_SST}, and AMSR-E (microwave) wind
    speed / water vapor / cloud liquid water / SST, plus
    \texttt{MODIS\_SST\_on\_AMSR\_E\_grid} and
    \texttt{convolved\_MODIS\_on\_AMSR\_E\_grid}. Reading this group
    alongside the root makes the file a MODIS\,$\leftrightarrow$\,AMSR-E
    co-located analysis product, not a vanilla L2 SST file.
  \item \texttt{/contributing\_granules/} --- provenance: per-orbit
    list of the source L2 granules that feed the orbit, with start
    and end times, byte ranges, and filenames. Pure metadata, no
    science.
\end{itemize}

Caveat to record for later: AMSR-E stopped producing data in October
2011, so the AMSR-E-named variables in a 2024 orbit file are presumably
all fill. The schema persists; the data should not. This needs
verification with a small slice on a 2024 vs.\ a 2010 file before any
claim is made in print.

\subsection{Metadata fetch: gradients\_by\_period sample}

Picked \texttt{augmented\_monthly\_stats\_07\_2020.nc} (July 2020).
DMR is 141\,lines total \emph{including} all attributes --- much
sparser than the orbit file. Dimensions \texttt{lon}=360,
\texttt{lat}=180. 34 variables, all Float64 of shape (lon, lat).
\textbf{Every variable in this dataset has zero attributes}. There are
\emph{no} global (root) attributes. No
\texttt{units}, no \texttt{long\_name}, no \texttt{standard\_name}, no
\texttt{\_FillValue}, no conventions tag. No coordinate variables
\texttt{lat} or \texttt{lon} are present --- only the dimensions of
those names.

This is the most striking finding so far and the most relevant to
\texttt{grep-dap}'s reason for being: a dataset where everything one
needs to know has to be inferred from \emph{structure plus names}
alone. The variable naming is regular enough to be tractable:

\begin{lstlisting}
day_pixel_count                 night_pixel_count
day_sum_SST                     night_sum_SST
day_sum_SST_squared             night_sum_SST_squared
day_as_pixel_count              night_as_pixel_count   # along-scan
day_sum_grad_as_per_km          night_sum_grad_as_per_km
day_sum_grad_as_per_km_squared  night_sum_grad_as_per_km_squared
day_at_pixel_count              night_at_pixel_count   # along-track
day_sum_grad_at_per_km          night_sum_grad_at_per_km
day_sum_grad_at_per_km_squared  night_sum_grad_at_per_km_squared
day_sum_grad_mag_per_km         night_sum_grad_mag_per_km
day_sum_grad_mag_per_km_squared night_sum_grad_mag_per_km_squared
day_sum_eastward_gradient       night_sum_eastward_gradient
day_sum_eastward_gradient_squared night_sum_eastward_gradient_squared
day_sum_northward_gradient      night_sum_northward_gradient
day_sum_northward_gradient_squared night_sum_northward_gradient_squared
day_sum_magnitude_gradient      night_sum_magnitude_gradient
day_sum_magnitude_gradient_squared night_sum_magnitude_gradient_squared
\end{lstlisting}

The structural inference, written out:
\begin{itemize}
  \item Each variable is a (lon, lat) grid of size $360\times180$.
    The dimension names alone are the strongest signal that this is
    a $1^\circ \times 1^\circ$ global grid. \emph{Caveat:} no
    coordinate variable names are present, so we cannot tell from
    the metadata whether \texttt{lat} runs $-90 \to +90$ or the
    reverse, nor whether \texttt{lon} is $0\to 360$ or $-180\to
    +180$.
  \item Day / night split is a strong hint that the upstream source
    is satellite, and the gradient and SST quantities tie directly
    to the variables in \texttt{SST\_Orbits/}. Day/night labels are
    a standard MODIS-SST partition (Aqua's $\sim$13:30 local
    ascending vs.\ $\sim$01:30 descending nodes).
  \item For each gradient flavor we get $(\Sigma, \Sigma^2, N)$, the
    canonical bookkeeping for computing means and variances after
    aggregation across periods, locations, or both.
  \item Three pixel-count flavors per period:
    \begin{itemize}
      \item \texttt{pixel\_count} --- valid SST pixels;
      \item \texttt{as\_pixel\_count} --- valid \emph{along-scan}
        gradient pixels (these can be invalid even when SST is valid,
        e.g.\ swath-edge effects);
      \item \texttt{at\_pixel\_count} --- valid \emph{along-track}
        gradient pixels.
    \end{itemize}
    The fact that there are separate $N$'s for SST, for along-scan
    gradients, and for along-track gradients confirms gradients are
    masked more aggressively than SST itself.
  \item The pair (\texttt{eastward\_gradient},
    \texttt{northward\_gradient}) and the pair
    (\texttt{grad\_as}, \texttt{grad\_at}) are two coordinate
    systems for the same physical vector field --- swath-relative
    (along/across scan) and geographic (east/north). The
    \texttt{magnitude\_gradient} and \texttt{grad\_mag\_per\_km} are
    the corresponding magnitudes. \emph{Caveat:} I have not yet
    verified that \texttt{sum\_grad\_mag\_per\_km} equals
    $\sqrt{\Sigma_{\rm east}^2 + \Sigma_{\rm north}^2}$ at the
    pixel level vs.\ the cell level; the bookkeeping per cell is
    almost certainly the sum of pixel magnitudes, not the magnitude
    of summed components.
\end{itemize}

\subsection{Data sampling}

Implemented in \texttt{scripts/uri\_sample\_data.py} and
\texttt{scripts/uri\_sample\_equator.py}. Two reads against the orbit
file (polar, equator) plus one full-grid read for the gradient file.
Highlights:

\begin{itemize}
  \item Orbit polar sample at $i_0=20000$ (nadir $\sim 78\text{--}79^\circ$N,
    Kara Sea region): \texttt{SST\_In} comes back $-15$ to $-3\,^\circ$C
    when the raw Int16 is scaled by 0.005. The declared
    \texttt{valid\_min/valid\_max} for \texttt{SST\_In} (raw) is
    $-600/9000$, i.e.\ $-3$ to $45\,^\circ$C; values below $-3\,^\circ$C
    are outside that valid range and should be treated as ice / land
    / unconverged. Both \texttt{eastward\_gradient} and
    \texttt{northward\_gradient} were entirely masked at this
    latitude (\texttt{n=0} after \texttt{\_FillValue} removal).
  \item Orbit equatorial sample at $i_0 \approx 29259$
    (nadir $\sim 0^\circ$, longitude $\sim -7^\circ$E, off the
    African coast): \texttt{SST\_In} in $[16.2, 29.0]\,^\circ$C,
    mean $27.4$, std $1.7$. Both gradient components populate at
    $\sim$20\,\% of pixels (1993 of 10000), with $|\nabla T|$
    samples in $\pm 0.5\,^\circ$C/km, mean near zero --- realistic
    for open tropical ocean.
  \item Lat/lon Int32 values divided by 1000 reproduce the
    Float32 \texttt{nadir\_latitude/longitude} to within $0.5^\circ$
    over the sample. Confirms the inferred implicit scale of
    $10^{-3}$ degree on the Int32 pixel coordinates.
  \item \texttt{regridded\_sst} (Int16, scale 0.005) returns values
    in $[26.8, 29.0]\,^\circ$C at the equator slice with fewer valid
    pixels (3662 vs.\ 10000 for raw SST\_In), consistent with the
    name suggesting a masked / re-gridded version.
  \item Gradient stats file (July 2020):
    \texttt{day\_pixel\_count} max $\sim$2.0$\times 10^5$ per
    $1^\circ$ cell, consistent with $\sim$14 day-orbits $\times$
    31 d $\times$ (pixel area / cell area). Per-cell mean SST
    derived as $\Sigma_{\rm SST}/N_{\rm pix}$ shows NaN propagation
    in some cells where \texttt{day\_sum\_SST} is NaN even though
    \texttt{day\_pixel\_count} is finite; this is a fill-value
    convention to resolve before any climatology is constructed
    (likely some cells are NaN by design where SST sums underflow
    or are flagged).
  \item Per-cell mean magnitude gradient $\sim 0.083\,^\circ$C/km,
    with a tail to $\sim 22\,^\circ$C/km. The tail is almost
    certainly numerical (e.g.\ coast / cloud-edge dropouts). The
    bulk is consistent with typical ocean SST gradients of
    $0.01\text{--}0.5\,^\circ$C/km.
\end{itemize}

\subsection{Things that bit me / would bite anyone}

\begin{itemize}
  \item Recursive catalog walks against this server are slow:
    individual \texttt{contents.html} requests take 0.5--2\,s, and
    the year/month structure of \texttt{SST\_Orbits/} (yearly
    pages are 20--30\,kB and the flat
    \texttt{gradients\_by\_period/} index is 2.9\,MB) means even a
    breadth-first walk at depth 2 with timeouts has to be sampled
    aggressively. My first walk script timed out at 45\,s after
    starting on the gradients directory; I refactored to walk the
    branches individually and to cap samples.
  \item Implicit scale on \texttt{latitude/longitude} (Int32 raw,
    1000$\times$ degrees) is not in the DMR. Any client that
    reads the orbit file naively gets garbage coordinates. This is
    exactly the class of bug \texttt{grep-dap} should flag: a
    declared \texttt{units=degrees\_north} on an Int32 with raw
    magnitudes near $10^5$ is internally inconsistent unless an
    implicit scale is assumed.
  \item The 403 directories are quietly informative even when
    blocked. The names alone partition into satellite-sensor
    sources (\texttt{JAXA\_Orbits}, \texttt{RSS\_Orbits},
    \texttt{MUR}), in-situ data (\texttt{iQuamBuoy}), matchup
    products (\texttt{matchups[\_v2,\_v3]}), and time-series
    products (\texttt{timeSeries[\_v2]}). \texttt{grep-dap} should
    not throw away the names of unreadable directories: they are
    part of the catalog signal.
\end{itemize}

\section{Process notes (carried forward)}

\begin{itemize}
  \item Variable / file names at URI's GSO Physical Oceanography
    group have historically referred to satellite SST products
    (PATHFINDER, AVHRR, GOES, MODIS-Aqua). The host name itself
    (\texttt{sst-aqua}) and the file naming
    (\texttt{AQUA\_MODIS\_orbit\_\dots\_L2\_SST-URI\_24-2.nc4})
    confirm MODIS-Aqua SST. Expectation and evidence agree here,
    but the project rule (\texttt{CLAUDE.md}: ``do not make up
    data'') still requires me to record both separately.
  \item For each claim about a variable I record whether it came
    from (a) a directly-read attribute, (b) a structural inference
    from the DMR, or (c) a statistical observation on a sampled
    slice. In the findings document the same convention is enforced
    with explicit Evidence tags.
  \item Scripts that produced these results are checked in under
    \texttt{scripts/}:
    \texttt{uri\_root\_walk.py},
    \texttt{uri\_walk\_public.py},
    \texttt{uri\_walk\_sst\_orbits.py},
    \texttt{uri\_sample\_data.py},
    \texttt{uri\_sample\_equator.py}.
\end{itemize}

\section{Session timing}

Started 2026-05-16 14:01\,UTC. Wrote final
\texttt{uri\_test\_case.tex} and closed the log at
2026-05-16 $\sim$14:25\,UTC. Within the one-hour budget.

\end{document}
